% ============================================================
%  00preambulo.tex  -  Preambulo consolidado del Plan de Tesis
%  (APA7 + espanol + graficos + matematica + matriz apaisada)
%  Compilar desde main.tex:  pdflatex -> biber -> pdflatex -> pdflatex
% ============================================================

% ---- Codificacion e idioma ----
\usepackage[utf8]{inputenc}
\usepackage[T1]{fontenc}
\usepackage{lmodern}
\usepackage[spanish,es-tabla]{babel}
\usepackage{csquotes}

% ---- Geometria (margenes tipo tesis) ----
\usepackage[a4paper,left=3cm,right=2.5cm,top=2.5cm,bottom=2.5cm]{geometry}

% ---- Interlineado 1.5 ----
\usepackage{setspace}
\onehalfspacing

% ---- Tipografia de titulos (se re-define mas abajo para APA7) ----
\usepackage{titlesec}

% ---- Tablas y figuras ----
\usepackage{graphicx}
\usepackage{booktabs}
\usepackage{array}
\usepackage{multirow}
\usepackage{tabularx}
\usepackage{float}
\usepackage[font=small,labelfont=bf]{caption}

% ---- Matematica (atencion, perdida, flujo de datos de la cascada) ----
\usepackage{amsmath}

% ---- Pagina apaisada para la matriz de consistencia ----
\usepackage{pdflscape}

% ---- Color de celdas / rampas ----
\usepackage{colortbl}

% ---- Graficos vectoriales (pgfplots / tikz) ----
\usepackage{pgfplots}
\pgfplotsset{compat=1.18}
\usepackage{tikz}
\usetikzlibrary{shapes.geometric,arrows.meta,positioning,calc}

% ---- Colores base ----
\usepackage{xcolor}
\definecolor{azulUNSAAC}{RGB}{0,71,133}
\definecolor{grisSuave}{RGB}{240,240,240}

% ---- Rutas de imagenes (logo, figuras) ----
\graphicspath{{imagenes/}{./}{../}}

% ---- Bibliografia APA7 (biblatex + biber) ----
\usepackage[
  backend=biber,
  style=apa,
  sortcites=true,
  sorting=nyt,
  language=spanish
]{biblatex}
\DeclareLanguageMapping{spanish}{spanish-apa}
% Fuentes: referencias generales (raiz del repo) y antecedentes de la revision.
\addbibresource{../referencias.bib}
\addbibresource{antecedentes.bib}

% ---- Listas e indices (tocloft ANTES de hyperref) ----
\usepackage{longtable}
\usepackage{algorithm}
\usepackage{algpseudocode}
\usepackage{tocloft}

% ---- Hipervinculos ----
\usepackage[hidelinks]{hyperref}
\hypersetup{
  colorlinks=true,
  linkcolor=black,
  citecolor=azulUNSAAC,
  urlcolor=azulUNSAAC
}

% ---- Icono ORCID ----
\usepackage{orcidlink}

% ---- Sangria de primer parrafo ----
\usepackage{indentfirst}

% ============================================================
%  Colores de la "ecuacion universal" del titulo (cafe/rojo/verde)
%  Esquema estandarizado (marco del docente): el titulo se descompone en tres
%  componentes, cada uno con un color fijo reutilizable en todo el proyecto.
% ============================================================
\definecolor{colProceso}{RGB}{122,66,20}   % cafe  -> Proceso / Accion
\definecolor{colMetodo}{RGB}{176,42,42}     % rojo  -> Tecnologia / Metodo
\definecolor{colContexto}{RGB}{31,110,54}   % verde -> Contexto / Aplicacion

% ============================================================
%  CORRECCIONES A OBSERVACIONES (revision del asesor)
% ============================================================
\definecolor{verdeOK}{RGB}{13,124,51}
% R1 consolidado: correcciones de la 1a ronda en NEGRO (texto normal).
\newcommand{\nuevo}[1]{#1}
\newenvironment{bloqueNuevo}{\par}{\par}
% R2 (2a ronda): aceptada; en NEGRO (texto normal).
\newcommand{\obsdos}[1]{#1}
\newenvironment{bloqueObsDos}{\par}{\par}
% R3: integrada al texto en NEGRO (versi{\'o}n para el jurado, sin resaltar correcciones).
\newcommand{\obstres}[1]{#1}
\newenvironment{bloqueObsTres}{\par}{\par}
% Mejoras propias de rigor (NO son una ronda del asesor): quedan en NEGRO (texto normal),
% integradas al documento como el resto. El macro se conserva por si luego se desea resaltar.
\newcommand{\obscuatro}[1]{#1}
\newenvironment{bloqueObsCuatro}{\par}{\par}

% ---- APA 7: numero y titulo ENCIMA de tablas y figuras; Nota al pie ----
\captionsetup{labelsep=newline,labelfont=bf,textfont=it,
              justification=raggedright,singlelinecheck=false}
\captionsetup[table]{position=above,skip=4pt}
\captionsetup[figure]{position=above,skip=6pt}
\newcommand{\notaAPA}[1]{\par\vspace{2pt}%
  \begin{minipage}{\linewidth}\footnotesize\textit{Nota.} #1\end{minipage}\par}

% ---- Estilo de tabla del proyecto (APA: cabecera negra sin relleno) ----
\newcommand{\hdr}[1]{\textbf{#1}}
\newcommand{\celnivel}[1]{\textbf{#1}}

% ---- Enlaces internos: siglas (quz/quy) y terminos del glosario ----
\DeclareRobustCommand{\sig}[1]{\hyperlink{sig:#1}{\texttt{#1}}}
\DeclareRobustCommand{\glo}[2]{\hyperlink{glo:#1}{#2}}
\newcommand{\destino}[1]{\hypertarget{#1}{}}

% ---- APA 7: encabezados ----
\titleformat{\section}[display]{\normalfont\Large\bfseries\filcenter}{Cap{\'i}tulo~\thesection}{0.3em}{}
\titleformat{name=\section,numberless}{\normalfont\Large\bfseries\filcenter}{}{0pt}{}
\titleformat{\subsection}{\normalfont\large\bfseries}{\thesubsection}{1em}{}
\titleformat{\subsubsection}{\normalfont\normalsize\bfseries\itshape}{\thesubsubsection}{1em}{}

% ---- Rampa estandar de los graficos: azul claro -> navy -> turquesa -> verde ----
\definecolor{g1}{HTML}{8AC4E6}
\definecolor{g2}{HTML}{66AEDA}
\definecolor{g3}{HTML}{4294CC}
\definecolor{g4}{HTML}{2E6DA0}
\definecolor{g5}{HTML}{1E4870}
\definecolor{g6}{HTML}{1F6476}
\definecolor{g7}{HTML}{229184}
\definecolor{g8}{HTML}{27A597}
\definecolor{g9}{HTML}{40B183}
\definecolor{g10}{HTML}{59BD6F}
\definecolor{g11}{HTML}{6AC362}
\definecolor{g12}{HTML}{7BC95A}

% ---- Numeracion por capitulos (I, II, ...) y subsecciones (2.1, 2.1.1) ----
\renewcommand{\thesection}{\Roman{section}}
\renewcommand{\thesubsection}{\arabic{section}.\arabic{subsection}}
\renewcommand{\thesubsubsection}{\arabic{section}.\arabic{subsection}.\arabic{subsubsection}}
\setcounter{secnumdepth}{3}
\setcounter{tocdepth}{3}

% ---- Nombres en espanol de los indices ----
\addto\captionsspanish{\renewcommand{\contentsname}{\'Indice general}}
\floatname{algorithm}{Algoritmo}
\renewcommand{\listalgorithmname}{\'Indice de algoritmos}

% ---- Indice de ecuaciones (lista propia via tocloft) ----
\newlistof{ecuaciones}{loe}{\'Indice de ecuaciones}
\newcommand{\regecuacion}[1]{%
  \addcontentsline{loe}{ecuaciones}{\protect\numberline{\theequation}#1}}
\let\cftecuacionesfont\cfttabfont
\let\cftecuacionespagefont\cfttabpagefont
\setlength{\cftecuacionesindent}{\cfttabindent}
\setlength{\cftecuacionesnumwidth}{\cfttabnumwidth}

% ---- Paginacion APA 7: numero en la esquina SUPERIOR DERECHA ----
\usepackage{fancyhdr}
\pagestyle{fancy}
\setlength{\headheight}{14.5pt}
\fancyhf{}
\fancyhead[R]{\thepage}
\renewcommand{\headrulewidth}{0pt}
\fancypagestyle{plain}{%
  \fancyhf{}%
  \fancyhead[R]{\thepage}%
  \renewcommand{\headrulewidth}{0pt}%
}
